\chapter{Formal Languages and the Chomsky Hierarchy}

Part~I developed analytic combinatorics in a setting-agnostic way: one starts with a class of objects, encodes it by a generating function, and studies the singularities of that function. This chapter connects that viewpoint to a classical source of combinatorial classes, namely \emph{formal languages}. The Chomsky hierarchy organizes those languages into four broad levels:
\begin{enumerate}
  \item regular,
  \item context-free,
  \item context-sensitive,
  \item recursively enumerable.
\end{enumerate}
The goal here is mostly orienting rather than proving every theorem from scratch. Results such as Myhill--Nerode, Kleene's theorem, the pumping lemmas, and the equivalence between context-free grammars and pushdown automata are major theorems of formal-language theory. I will use them as imported facts and focus on the analytic-combinatorial lessons they support.

The main lesson is one-way, not one-to-one: regular languages always lead to rational length generating functions, and unambiguous context-free languages lead to algebraic ones. But moving higher in the hierarchy does \emph{not} create a perfect analytic classification. Higher-level languages can still have simple generating functions when counted by length. So the hierarchy is a guide to tractability, not a complete dictionary from language class to analytic type.

\section{Alphabets, Strings, and Languages}

\begin{definition}
An \emph{alphabet} $\Sigma$ is a finite nonempty set. A \emph{string} over $\Sigma$ is a finite sequence
\[
  w = \sigma_1\sigma_2\cdots\sigma_k,
  \qquad \sigma_i \in \Sigma.
\]
Its \emph{length} is $|w|=k$. The unique string of length $0$ is the \emph{empty string} $\varepsilon$.
\end{definition}

For example, over $\Sigma=\{a,b\}$, the word `abba' has length $4$.

The set $\Sigma^*$ of all finite strings over $\Sigma$ is called the \emph{free monoid} on $\Sigma$; this simply means that strings can be concatenated and that $\varepsilon$ acts as the identity element. We write
\[
  \Sigma^n = \{w \in \Sigma^* : |w|=n\},
  \qquad
  \Sigma^+ = \Sigma^* \setminus \{\varepsilon\}.
\]
A \emph{language} is any subset $L \subseteq \Sigma^*$.

The counting object relevant for this book is
\[
  L(z) = \sum_{n \ge 0} |L \cap \Sigma^n| z^n,
\]
which records the number of accepted words of each length. For the full language $\Sigma^*$ over a two-letter alphabet,
\[
  |\Sigma^n| = 2^n,
  \qquad
  \sum_{n \ge 0} 2^n z^n = \frac{1}{1-2z}.
\]

\section{Finite Automata and Regular Languages}

\begin{definition}
A \emph{deterministic finite automaton} (DFA) is a tuple
\[
  M = (Q,\Sigma,\delta,q_0,F)
\]
where $Q$ is a finite set of states, $\Sigma$ is the input alphabet, $\delta : Q \times \Sigma \to Q$ is the transition function, $q_0 \in Q$ is the start state, and $F \subseteq Q$ is the set of accepting states.
\end{definition}

The transition function extends from single letters to whole strings by the recursion
\[
  \hat\delta(q,\varepsilon)=q,
  \qquad
  \hat\delta(q,w\sigma)=\delta(\hat\delta(q,w),\sigma).
\]
So to run a DFA on a word, one simply feeds the letters in one at a time. For instance, in the parity automaton below, the word `abaa' toggles parity three times and ends in the odd state.

\begin{example}[Even number of $a$'s]
Let $\Sigma=\{a,b\}$. The language
\[
  L = \{w : w \text{ contains an even number of } a\text{'s}\}
\]
is recognized by a DFA with two states: ``even'' and ``odd.'' Reading an $a$ switches the state, while reading a $b$ leaves it unchanged. The start state is ``even,'' and it is also the only accepting state.

\begin{center}
\begin{tikzpicture}[shorten >=1pt, node distance=2.5cm, on grid, auto,
    every state/.style={circle, draw, minimum size=1cm, text centered}]
  
  \node[state, accepting, initial, initial text={}] (q0) {even};
  \node[state] (q1) [right=of q0] {odd};
  
  \path[->] 
    (q0) edge [bend left] node [above] {$a$} (q1)
         edge [loop above] node {$b$} ()
    (q1) edge [bend left] node [below] {$a$} (q0)
         edge [loop above] node {$b$} ();
\end{tikzpicture}
\end{center}
\end{example}

A \emph{nondeterministic finite automaton} (NFA) allows several possible next states, and usually also $\varepsilon$-transitions. One may therefore think of its transition map as taking values in subsets of states over letters in $\Sigma \cup \{\varepsilon\}$. The key determinization theorem is this:

\begin{theorem}[NFA $=$ DFA, used as a black box]
Every NFA accepts the same language as some DFA.
\end{theorem}

The idea is the subset construction: a DFA state remembers the whole set of NFA states that are currently possible after reading the input so far.

\begin{definition}
A language is \emph{regular} if it is accepted by some DFA (equivalently, by some NFA).
\end{definition}

\section{Myhill--Nerode and Regularity}

For a language $L$, define an equivalence relation on strings by
\[
  x \sim_L y
  \quad\Longleftrightarrow\quad
  \forall z \in \Sigma^*\; (xz \in L \iff yz \in L).
\]
So two strings are equivalent if they have exactly the same possible futures leading to acceptance.

\begin{theorem}[Myhill--Nerode, used as a black box]
A language $L$ is regular if and only if $\sim_L$ has only finitely many equivalence classes. In that case, the number of classes is exactly the number of states in the minimal DFA for $L$.
\end{theorem}

The forward direction matches intuition: if two strings lead a DFA to the same state, then every future continuation is processed identically. The converse constructs a DFA whose states are the equivalence classes themselves.

This theorem is a powerful nonregularity test. For
\[
  L = \{a^n b^n : n \ge 0\},
\]
the strings
\[
  \varepsilon, a, a^2, a^3, \dots
\]
are pairwise inequivalent, because
\[
  a^i b^i \in L,
  \qquad
  a^j b^i \notin L \quad (j \ne i).
\]
So $\sim_L$ has infinitely many classes, and the language is not regular.

\section{Regular Expressions}

Regular expressions give a syntactic description of regular languages. Their semantics are:
\begin{itemize}
  \item $\emptyset$ denotes the empty language,
  \item $\varepsilon$ denotes the language $\{\varepsilon\}$,
  \item a symbol $\sigma$ denotes the one-word language $\{\sigma\}$,
  \item $R+S$ denotes union,
  \item $RS$ denotes concatenation:
  \[
    L(RS) = \{xy : x \in L(R),\; y \in L(S)\},
  \]
  \item $R^*$ denotes Kleene star:
  \[
    L(R^*) = \bigcup_{k \ge 0} L(R)^k.
  \]
\end{itemize}

\begin{theorem}[Kleene, used as a black box]
A language is regular if and only if it is denoted by some regular expression.
\end{theorem}

There is an important warning for this book. Regular-expression operations do \emph{not} automatically translate into symbolic-method OGF operations when we count \emph{distinct words}. Union may fail to be disjoint, concatenation may fail to be uniquely decodable, and star may be ambiguous. So the analogy between regexes and symbolic constructors is only safe under uniqueness / disjointness hypotheses. For counting distinct accepted words by length, automata-based linear algebra is often the safer route.

\section{The Pumping Lemma for Regular Languages}

\begin{theorem}[Pumping lemma for regular languages, used as a black box]
If $L$ is regular, then there exists a pumping length $p$ such that every word $w \in L$ with $|w| \ge p$ can be written as
\[
  w = xyz
\]
with $|xy| \le p$, $|y| \ge 1$, and
\[
  xy^k z \in L \qquad \text{for all } k \ge 0.
\]
\end{theorem}

The proof uses repeated states along the run of a DFA. It is important to remember that the pumping lemma is only a \emph{necessary} condition for regularity, not a sufficient one.

For $\{a^n b^n : n \ge 0\}$, take $w=a^p b^p$. Since the first $p$ symbols are all $a$, the pumped block $y$ lies entirely inside the initial block of $a$'s. Pumping $y$ changes the number of $a$'s without changing the number of $b$'s, so the pumped word leaves the language.

\section{Context-Free Grammars and Pushdown Automata}

\begin{definition}
A \emph{context-free grammar} (CFG) is a tuple
\[
  G = (V,\Sigma,P,S)
\]
where $V$ is a finite set of nonterminals, $\Sigma$ is a finite set of terminals, $P$ is a finite set of productions
\[
  A \to \alpha,
  \qquad A \in V,
  \quad \alpha \in (V \cup \Sigma)^*,
\]
and $S \in V$ is the start symbol.
\end{definition}

A language is \emph{context-free} if it is generated by some CFG.

\begin{example}[$a^n b^n$]
The grammar
\[
  S \to aSb \mid \varepsilon
\]
generates exactly $\{a^n b^n : n \ge 0\}$.
\end{example}

\begin{example}[Dyck language]
Balanced parenthesis strings are generated unambiguously by the first-return grammar
\[
  S \to (S)S \mid \varepsilon.
\]
If $z$ counts individual parentheses, this gives
\[
  S(z) = 1 + z^2 S(z)^2.
\]
So the generating function is
\[
  S(z) = \frac{1-\sqrt{1-4z^2}}{2z^2},
\]
which is algebraic rather than rational.
\end{example}

A \emph{pushdown automaton} (PDA) is a finite automaton equipped with a stack. The stack records pending nested structure, which is exactly what recursive grammatical structure needs.

\begin{theorem}[CFG $=$ PDA, used as a black box]
A language is context-free if and only if it is accepted by some pushdown automaton.
\end{theorem}

The intuition is that grammar expansion creates obligations that must be completed later, and the stack stores those obligations.

\section{The Pumping Lemma for Context-Free Languages}

\begin{theorem}[Pumping lemma for context-free languages, used as a black box]
If $L$ is context-free, then there exists a pumping length $p$ such that every sufficiently long word in $L$ can be written as
\[
  w = uvxyz
\]
with
\[
  |vxy| \le p,
  \qquad
  |vy| \ge 1,
  \qquad
  uv^k x y^k z \in L \text{ for all } k \ge 0.
\]
\end{theorem}

Standard proofs use parse trees in a controlled normal form. As with the regular pumping lemma, this is only a necessary condition for context-freeness, not a sufficient one.

To show that
\[
  \{a^n b^n c^n : n \ge 0\}
\]
is not context-free, one studies a long word $a^p b^p c^p$. The short middle region $vxy$ can touch at most two of the three blocks, so pumping it inevitably changes one or two symbol counts while leaving the third fixed. Equal counts cannot be preserved.

\section{The Chomsky Hierarchy}

The hierarchy is best read as a sequence of increasingly powerful language models:
\begin{itemize}
  \item \textbf{Type 3: regular languages.} Recognized by finite automata. Their memory is finite.
  \item \textbf{Type 2: context-free languages.} Recognized by pushdown automata. They can handle one stack of nested dependencies.
  \item \textbf{Type 1: context-sensitive languages.} Recognized by linear-bounded automata.
  \item \textbf{Type 0: recursively enumerable languages.} Recognized by unrestricted Turing machines.
\end{itemize}
The inclusions are strict.

For example, regular languages are closed under union, intersection, complement, concatenation, and star. Context-free languages are closed under union, concatenation, and star, but not in general under intersection or complement. A standard witness is
\[
  \{a^n b^n c^k : n,k \ge 0\} \cap \{a^k b^n c^n : k,n \ge 0\} = \{a^n b^n c^n : n \ge 0\},
\]
and the right-hand side is not context-free.

We will not dwell on the upper two levels in detail. For this book, they mainly play a cautionary role: once one leaves the regular and unambiguous context-free worlds, the clean analytic methods of Parts~I--II become less automatic.

\section{The Analytic-Combinatorial Payoff}

The cleanest implications for this book are these.

\begin{enumerate}
  \item If $L$ is regular, then its length generating function is rational.
  \item If $L$ is generated by an unambiguous CFG, then its length generating function is algebraic.
\end{enumerate}

The first claim can be seen by linear algebra. Let a DFA have states $q_1,\dots,q_m$, start-state vector $u$, and accepting-state vector $v$. Define a matrix $A$ by letting $A_{ij}$ be the number of alphabet symbols that send state $q_i$ to state $q_j$. Then the number of accepted words of length $n$ is
\[
  u^T A^n v,
\]
so the OGF is
\[
  L(z) = \sum_{n \ge 0} u^T A^n v\, z^n = u^T (I-zA)^{-1} v,
\]
which is rational. The $n=0$ term automatically handles the empty word: it contributes $u^T v$, which is $1$ exactly when the start state is accepting.

The second claim is the Chomsky--Sch\"utzenberger theorem from Chapter~5. Here again, the direction is one-way and structural: unambiguous CFGs give algebraic generating functions, but algebraicity alone does not force a unique asymptotic regime without the additional hypotheses discussed in Chapter~5.

What should one say about higher levels? Not that ``context-sensitive'' means one specific analytic type. The honest statement is weaker: once we leave the regular and unambiguous context-free settings, we lose the general theorems guaranteeing rational or algebraic generating functions. Some higher-level languages may still have rational or algebraic length generating functions, but there is no comparably clean universal structure theorem waiting for us there.

One bound remains universal. For every language $L \subseteq \Sigma^*$ over a finite alphabet,
\[
  |L \cap \Sigma^n| \le |\Sigma|^n.
\]
So word counts by length are always geometrically bounded. The ``exotic'' behavior that can arise at higher levels is about complicated singularity structure or the failure of easy continuation theorems, not about supergeometric growth of the counting sequence.

This is why regular and unambiguous context-free languages are analytically privileged: they come with robust structural theorems that turn formal-language descriptions into generating functions we can actually study.

\section*{Preview of Chapter 7}

Chapter~7 studies \emph{weighted finite automata}, a probabilistic refinement of regular-language automata. There the rational generating-function viewpoint becomes computational: matrices, resolvents, and spectral radii will let us count, weight, and eventually approximate sequence models.

\section*{Further Reading}

\begin{itemize}
  \item \textbf{J. E. Hopcroft and J. D. Ullman}, \emph{Introduction to Automata Theory, Languages, and Computation} (1979) \cite{HopcroftUllman1979} --- the classical reference for automata, grammars, Myhill--Nerode, pumping lemmas, and the Chomsky hierarchy.
  \item \textbf{M. Sipser}, \emph{Introduction to the Theory of Computation} (2013) \cite{Sipser2013} --- a modern, readable introduction to the same material.
  \item \textbf{P. Flajolet and R. Sedgewick}, \emph{Analytic Combinatorics} (2009) \cite{FlajoletSedgewick2009} --- for the symbolic and generating-function side of the regular / context-free story.
  \item \textbf{J. Berstel and C. Reutenauer}, \emph{Noncommutative Rational Series with Applications} (2011) \cite{BerstelReutenauer2011} --- for the algebraic background behind rational series and weighted automata.
\end{itemize}
